\documentclass[10pt, letterpaper]{article}

% Packages:
\usepackage[
    ignoreheadfoot,
    top=2 cm,
    bottom=2 cm,
    left=2 cm,
    right=2 cm,
    footskip=1.0 cm,
]{geometry}
\usepackage[explicit]{titlesec}
\usepackage{tabularx}
\usepackage{array}
\usepackage[dvipsnames]{xcolor}
\definecolor{primaryColor}{RGB}{0, 79, 144} % blue
\usepackage{enumitem}
\usepackage{fontawesome5}
\usepackage{amsmath}
\usepackage[
    pdftitle={Chunxu Han's CV},
    pdfauthor={Chunxu Han},
    pdfcreator={LaTeX with RenderCV},
    colorlinks=true,
    urlcolor=primaryColor
]{hyperref}
\usepackage[pscoord]{eso-pic}
\usepackage{calc}
\usepackage{bookmark}
\usepackage{lastpage}
\usepackage{changepage}
\usepackage{paracol}
\usepackage{ifthen}
\usepackage{needspace}
\usepackage{iftex}
\usepackage{multicol}
\usepackage{datetime2}
\usepackage{graphicx}

\ifPDFTeX
    \input{glyphtounicode}
    \pdfgentounicode=1
    \usepackage[T1]{fontenc}
    \usepackage[utf8]{inputenc}
    \usepackage{lmodern}
\fi

\usepackage[default, type1]{sourcesanspro}

\AtBeginEnvironment{adjustwidth}{\partopsep0pt}
\pagestyle{empty}
\setcounter{secnumdepth}{0}
\setlength{\parindent}{0pt}
\setlength{\topskip}{0pt}
\setlength{\columnsep}{0.15cm}
\makeatletter
\let\ps@customFooterStyle\ps@plain
\patchcmd{\ps@customFooterStyle}{\thepage}{
    \color{gray}\textit{\small Chunxu Han - Page \thepage{} of \pageref*{LastPage}}
}{}{}
\makeatother
\pagestyle{customFooterStyle}

\titleformat{\section}{
    \needspace{4\baselineskip}
    \Large\color{primaryColor}
}{}{}{
    \textbf{#1}\hspace{0.15cm}\titlerule[0.8pt]\hspace{-0.1cm}
}[]

\titlespacing{\section}{-1pt}{0.3 cm}{0.2 cm}

\newenvironment{highlights}{
    \begin{itemize}[
        topsep=0.10 cm,
        parsep=0.10 cm,
        partopsep=0pt,
        itemsep=0pt,
        leftmargin=0.4 cm + 10pt
    ]
}{
    \end{itemize}
}

\newenvironment{highlightsforbulletentries}{
    \begin{itemize}[
        topsep=0.10 cm,
        parsep=0.10 cm,
        partopsep=0pt,
        itemsep=0pt,
        leftmargin=10pt
    ]
}{
    \end{itemize}
}

\newenvironment{onecolentry}{
    \begin{adjustwidth}{0.2 cm + 0.00001 cm}{0.2 cm + 0.00001 cm}
}{
    \end{adjustwidth}
}

\newenvironment{twocolentry}[2][]{
    \onecolentry
    \def\secondColumn{#2}
    \setcolumnwidth{\fill, 4.5 cm}
    \begin{paracol}{2}
}{
    \switchcolumn \raggedleft \secondColumn
    \end{paracol}
    \endonecolentry
}

\newenvironment{threecolentry}[3][]{
    \onecolentry
    \def\thirdColumn{#3}
    \setcolumnwidth{1 cm, \fill, 4.5 cm}
    \begin{paracol}{3}
    {\raggedright #2} \switchcolumn
}{
    \switchcolumn \raggedleft \thirdColumn
    \end{paracol}
    \endonecolentry
}

\newenvironment{header}{
    \setlength{\topsep}{0pt}\par\kern\topsep\centering\color{primaryColor}\linespread{1.5}
}{
    \par\kern\topsep
}

\let\hrefWithoutArrow\href
\renewcommand{\href}[2]{\hrefWithoutArrow{#1}{\ifthenelse{\equal{#2}{}}{ }{#2 }\raisebox{.15ex}{\footnotesize \faExternalLink*}}}


\begin{document}
    \newcommand{\AND}{\unskip
        \cleaders\copy\ANDbox\hskip\wd\ANDbox
        \ignorespaces
    }
    \newsavebox\ANDbox
    \sbox\ANDbox{}

    \begin{header}
        \begin{minipage}[t]{0.78\linewidth}
        \raggedright
        \fontsize{28 pt}{28 pt}
        \textbf{Chunxu Han}

        \vspace{0.3 cm}

        \normalsize
        \mbox{{\footnotesize\faMapMarker*}\hspace*{0.13cm}Peder Skrams Gade 16A, København K, Denmark}%
        \kern 0.25 cm%
        \AND%
        \kern 0.25 cm%
        \mbox{\hrefWithoutArrow{mailto:s220311@dtu.dk}{{\footnotesize\faEnvelope[regular]}\hspace*{0.13cm}s220311@dtu.dk}}%
        \kern 0.25 cm%
        \AND%
        \kern 0.25 cm%
        \mbox{\hrefWithoutArrow{tel:+45 52 73 06 85}{{\footnotesize\faPhone*}\hspace*{0.13cm}+45 52 73 06 85}}%
        \kern 0.25 cm%
        \AND%
        \kern 0.25 cm%
        \mbox{\hrefWithoutArrow{https://www.linkedin.com/in/chunxu-han/}{{\footnotesize\faLinkedinIn}\hspace*{0.13cm}chunxu-han}}%
        \kern 0.25 cm%
        \AND%
        \kern 0.25 cm%
        \mbox{\hrefWithoutArrow{https://github.com/chunxubioinfor}{{\footnotesize\faGithub}\hspace*{0.13cm}chunxubioinfor}}%
        \end{minipage}%
        \hfill%
        \begin{minipage}[t]{0.18\linewidth}
        \raggedleft
        \vspace{-1.0cm}
        \includegraphics[width=2.4cm, height=2.8cm, keepaspectratio=false]{portrait.jpg}
        \end{minipage}
    \end{header}

    \vspace{0.4 cm - 0.3 cm}


    %% ===== TAILOR THIS: Relevant Experience (2–3 bold-lead paragraphs) =====
    %% Check source_materials/relevant_experience_bank.md first — reuse or adapt existing paragraphs.
    \section{Relevant Experience}

\begin{onecolentry}

\textbf{Quantitative Research \& Signal Development} \\
At Novo Nordisk, my thesis centred on developing a novel quantitative metric --- the Batch Mixing Index (BMI) --- to detect and measure systematic variation in large single-cell RNA-seq datasets. I designed the mathematical formulation, implemented it in Python, and validated it using permutation-based significance testing across multiple experimental conditions. This is the same loop as systematic strategy research: define a signal, build a pipeline to compute it, and test whether it holds up statistically.

        \vspace{0.3 cm}

\textbf{Large-Scale Data Pipelines \& Research Software} \\
At DTU HealthTech, I built data infrastructure for isoform-level proteomics analysis, curating and re-quantifying more than 40,000 samples from a public database using HPC-based parallel processing pipelines. I co-developed \href{https://github.com/kvittingseerup/IsoformSwitchAnalyzeR}{IsoformSwitchAnalyzeR} v2, an R package used by researchers worldwide, now published as a preprint on bioRxiv. For my thesis I built \href{https://github.com/chunxubioinfor/scbatch_vari}{scBatch}, a Python package for structured batch variability assessment. Writing clean, testable research software that others rely on is something I take seriously.

        \vspace{0.3 cm}

\textbf{Moving Between Theory and Implementation} \\
I adapt quickly across problem domains and enjoy the full cycle from mathematical formulation to working code. I transitioned from biotechnology to consulting to bioinformatics, and in each case built the technical skills needed from scratch. I hold DTU's \textbf{Machine Learning} and \textbf{Deep Learning} courses and have applied these methods in real research settings. I am comfortable in environments where ideas are judged by evidence and the bar for implementation quality is high.

\end{onecolentry}

    \section{Education}

        \begin{twocolentry}{
    Lyngby, Denmark

    Sept 2023 -- March 2026
}
    \textbf{MSc in Bioinformatics and Systems Biology}, Technical University of Denmark
    \begin{highlights}
        \item I enrolled in both \textbf{Machine Learning} and \textbf{Deep Learning} courses, enabling me to develop an \textbf{AI-driven mindset} when formulating projects and developing new bioinformatics pipelines or tools.
    \end{highlights}
\end{twocolentry}

        \vspace{0.3 cm}

\begin{twocolentry}{
    Beijing, China

    Sept 2017 -- June 2021
}
    \textbf{BSc in Biotechnology}, Beijing Normal University
    \begin{highlights}
        \item \textbf{Coursework:} Built a solid foundation in \textbf{Biotech} with courses covering ecology, molecular biology, cell biology, and bioinformatics, alongside laboratory training.
    \end{highlights}
\end{twocolentry}


    \section{Work Experience}

        \begin{twocolentry}{
    Målev, Denmark

    Sep 2025 -- March 2026
}
    \textbf{Master Student}, Novo Nordisk

    \vspace{0.1cm}

    My master's thesis focuses on developing computational frameworks to evaluate batch-to-batch variability in single-cell RNA-seq data:

    \begin{highlights}
        \item \textbf{Literature review and benchmark:} Systematically identify and evaluate existing computational tools for assessing batch comparability in both cell-type composition and transcriptional profiles

        \vspace{0.1cm}

        \item \textbf{Methodology development:} Design, develop and test quantitative metrics and simplified scoring strategies on in-house data to support batch evaluation in GMP-adjacent research settings.

        \vspace{0.1cm}

        \item \textbf{AI/ML approaches:} Explore machine learning methods to correlate transcriptomic variability with functional outcomes, contributing to predictive quality assessment.

        \vspace{0.1cm}

        \item \textbf{Meetings and Collaboration:} Participate in team meetings to stay updated on current trends in the Pharma industry. Collaborate with cross-functional R\&D and CMC teams to ensure the biological interpretability and industrial relevance.
    \end{highlights}
\end{twocolentry}

        \vspace{0.3 cm}

        \begin{twocolentry}{
    Lyngby, Denmark

    Oct 2023 -- July 2025
}
    \textbf{Bioinformatics Research Assistant}, DTU HealthTech

    \vspace{0.1cm}

    I worked in Isoform Analysis Group for two years, which has been a transformative experience. I have grown from a student who only focuses on theoretical studies into a bioinformatician to deal with real-world problems:

    \begin{highlights}
        \item \textbf{Development and Maintenance of Tools:} Contributed to developing and updating the \href{https://github.com/kvittingseerup/IsoformSwitchAnalyzeR}{IsoformSwitchAnalyzeR} R package. This experience sharpened my coding skills and can-do mindset---especially when handling issues with tools I'd never used before. I pushed myself to attempt and now I'm comfortable with the steep learning curve.

        \vspace{0.1cm}

        \item \textbf{Bioinformatics Pipelines:} Developed and optimized bioinformatics pipelines, supporting others in running analyses. I mastered multi-tasking, learned to leverage \textbf{HPC} for parallel processing, and \textbf{Git} to achieve version control.

        \vspace{0.1cm}

        \item \textbf{Ad-hoc Work:} Managed various ad-hoc tasks, from literature reviews to data organization, learning how to \textbf{prioritize and organize} work effectively in a fast-paced research environment.
    \end{highlights}
\end{twocolentry}

        \vspace{0.3 cm}

        \begin{twocolentry}{
    Shanghai, China

    Feb 2023 -- May 2023
}
    \textbf{Consulting Intern}, Illuminera, IQVIA

    \vspace{0.1cm}

    My first experience in a business setting, where I gained insights into the pharmaceutical market and developed professional skills:

    \begin{highlights}
        \item \textbf{Qualitative and Quantitative Research:} Conducted market trend and competitor analysis to formulate actionable strategies to clients. I also gained a deeper understanding of the pharmaceutical market in China and globally.

        \vspace{0.1cm}

        \item \textbf{Client and KOLs Engagement:} I learnt how to build professional relationships with clients and KOLs and coordinate with stakeholders to align project progress.

        \vspace{0.1cm}

        \item \textbf{Data-Driven Deliverables:} Utilized tools like PowerPoint and Power BI to deliver strategic reports, transforming complex data into actionable business insights.
    \end{highlights}
\end{twocolentry}

        \vspace{0.3 cm}


    \section{Hobby Projects}

\begin{onecolentry}

\textbf{\href{https://github.com/chunxubioinfor/DailyJobMatch}{DailyJobMatch} --- AI Job Matching Agent} \\
Built an AI agent that automatically fetches, filters, and ranks job postings daily based on my CV and preferences. The agent uses LLM-based scoring and sends a daily digest. This project taught me that AI becomes most useful when it is customised and built with real constraints in mind, rather than used as a generic chatbot.

        \vspace{0.3 cm}

\textbf{\href{https://github.com/chunxubioinfor/job-tracker}{Job Tracker} --- Full-Stack Web App} \\
Built a personal job application tracker using Next.js, Supabase, and Vercel. Features include a private applications pipeline with status tracking, CV/cover letter file storage, and a weekly reflection log. Designed with Row Level Security so multiple users can use it independently.

\end{onecolentry}

    \section{Skills and Hobbies}

        \begin{onecolentry}
            \textbf{Programming:} Python, R, Bash, SQL, Snakemake
        \end{onecolentry}

        \vspace{0.2 cm}

        \begin{onecolentry}
            \textbf{Languages:} English (proficient), Chinese (native), Danish (Basic)
        \end{onecolentry}

        \vspace{0.2 cm}

        \begin{onecolentry}
            \textbf{Soft Skills:} Communication, Teamwork, Client Engagement, Multitasking, Critical Thinking, Strong Work Ethic
        \end{onecolentry}

        \vspace{0.2 cm}

        \begin{onecolentry}
            \textbf{Hobbies:}
            I've been playing \textbf{badminton} since an early age---it's a hobby that keeps me focused and energized. During holidays, I enjoy \textbf{traveling} around Europe to experience new food and cultures. I also love \textbf{hiking}, both for the chance to connect with nature and for the challenge of seeing hidden scenery that isn't easily accessible. After moving to Denmark, I discovered a new passion: \textbf{sailing}. Last summer, we sailed to Ven and Hallands Väderö, and it was a completely new experience and I feel free to ride the infinite waves.
        \end{onecolentry}


\end{document}
